% ============================================================
%  TERM PROJECT – PROPOSAL (template)
%  Databases · THGA Bochum
%
%  How to use:
%    1. Copy this folder into your own repository.
%    2. Replace every  <<< ... >>>  placeholder with your content.
%    3. Build with  `make`  (see the Makefile / README).
%    4. Send the resulting PDF to the lecturer and wait for approval.
% ============================================================
\documentclass[11pt,a4paper]{article}

\usepackage{thga-db}

% ----- Fill in your metadata -------------------------------
\renewcommand{\thgaKurs}{Databases}
\renewcommand{\thgaDatum}{Summer Term 2026}
\newcommand{\authorName}{<<< Your name >>>}
\newcommand{\projectTitle}{<<< Working title of your application >>>}

\begin{document}
\thispagestyle{empty}
\thgaTitel{Term Project – Proposal}{\projectTitle}

\begin{infobox}
\textbf{Author:} \authorName \\
\textbf{Module:} Introduction to Database Management Systems · THGA Bochum \\
\textbf{Submitted to:} Stephan Bökelmann · \href{mailto:sboekelmann@ep1.rub.de}{sboekelmann@ep1.rub.de} \\
\textbf{Target size:} 2--4 pages
\end{infobox}

% ============================================================
\section{Application domain and motivation}
% ============================================================

<<< Which application do you build and which problem does it solve? Who uses
it, and for what? Add one or two concrete usage scenarios. >>>

% ============================================================
\section{Data model -- ER sketch}
% ============================================================

<<< Insert your ER sketch (a photographed hand drawing is fine). It must show
entity types, key attributes, relationships and cardinalities, including at
least one N:M relationship. >>>

% Example for including an image:
% \begin{center}
%   \includegraphics[width=0.9\linewidth]{er-sketch.jpg}
% \end{center}

% ============================================================
\section{From model to schema}
% ============================================================

<<< Derive the relational schema. For each table give the columns, the primary
key (underlined) and the foreign keys. >>>

\begin{itemize}
  \item \texttt{<<<table>>>}(\underline{id}, <<<column>>>, \dots)
  \item \texttt{<<<table>>>}(\underline{id}, \emph{<<<foreign\_key>>>}, \dots)
\end{itemize}

\textbf{Normalisation.} <<< One or two sentences: why is the schema free of
redundancy (at least 3NF)? >>>

% ============================================================
\section{API design}
% ============================================================

<<< List the planned FastAPI endpoints. Mark where a JOIN / an aggregation
sits, and which endpoints require the X-API-Key. >>>

\renewcommand{\arraystretch}{1.2}
\begin{tabularx}{\linewidth}{@{}p{2cm}p{4cm}X@{}}
\toprule
\textbf{Method} & \textbf{Path} & \textbf{Purpose} \\
\midrule
\texttt{GET}  & \texttt{/<<<…>>>} & <<< reads … (JOIN / aggregation) >>> \\
\texttt{POST} & \texttt{/<<<…>>>} & <<< creates … (requires X-API-Key) >>> \\
\bottomrule
\end{tabularx}

% ============================================================
\section{Frontend idea}
% ============================================================

<<< What does the frontend look like and which endpoints does it operate?
Python/tkinter is recommended; if you use something else, say what and why.
The frontend is delivered as a .deb installer and asks for the API URL and
X-API-Key in a connection dialog at startup. >>>

% ============================================================
\section{Operation with Docker Compose}
% ============================================================

<<< Which services run in your docker-compose.yml (at least postgres and api)?
Where do the credentials live (.env)? How is the backend deployed on the
lecture server? >>>

% ============================================================
\section{Scope, milestones and risks}
% ============================================================

\textbf{Scope.} <<< How many tables and endpoints? >>>

\textbf{Milestones.}
\begin{enumerate}[topsep=2pt]
  \item <<< e.g. schema + seed data ready >>>
  \item <<< e.g. API reading/writing >>>
  \item <<< e.g. frontend operates the core endpoints, packaged as .deb >>>
  \item <<< documentation + video >>>
\end{enumerate}

\textbf{Risks.} <<< One or two risks that could make the project fail. >>>

\vspace{1em}
\begin{infobox}
\textbf{Effort budget:} the whole term project (system, documentation and video)
should fit into roughly \textbf{40\,h of work} over at most \textbf{2 months}.
Keep your scope within that budget.
\end{infobox}

\end{document}
